\chapter{Introduction}
\label{ch:introduction}

\section{Research Background}

Modern game agents need to execute complex behavioural logic. For a small prototype, writing all decisions in a single procedural script is feasible; however, as states and interruption conditions increase, the control flow becomes increasingly difficult to inspect and modify. A behaviour tree (BT) represents decision logic as a hierarchy composed of control nodes and leaf nodes. Composite nodes determine how children are selected, while conditions and actions connect the abstract decision structure to game state and Actor code. High-level policies can be edited independently of movement, animation or combat methods, which is one of the principal values of behaviour trees in game development \parencite{colledanchise2018bt}.

Traditional node--link trees can directly express parent--child topology, but rapidly consume space as their width and depth grow. Behaviour-tree cards may also display parameters, descriptions, Decorators, blackboard keys and execution states, making them larger than ordinary hierarchies that display only labels. When dozens or hundreds of cards are placed on a limited canvas, developers must frequently zoom, pan and search, and the current execution path may be submerged among inactive branches. This is not merely an aesthetic issue, because the node graph is the main tool through which a designer locates tasks, understands order, modifies conditions and diagnoses failures.

This project therefore uses a functionally complete Godot behaviour-tree plugin as a platform for research into large-tree display. While preserving upper and lower ports and the semantic rule that siblings execute from left to right, it adds independently switchable focus+context, complexity-management, navigation and runtime-explanation techniques. Each technique can be disabled independently, both to prevent a malfunctioning visual feature from destabilising the editor and to retain a reproducible baseline comparison condition.

\section{Problem Definition and Research Gap}

The fields of tree and graph visualisation have proposed many methods. Fisheye and focus+context views allocate more space to elements near the focus \parencite{furnas1986fisheye,lamping1995hyperbolic}; overview+detail preserves a global reference view \parencite{cockburn2008review}; multi-column layouts can improve browsing in trees with large fan-outs \parencite{song2010largefanout}; incremental expansion and collapse can reduce graph complexity while attempting to preserve the user's mental map \parencite{dogrusoz2018complexity,misue1995mentalmap}; and research into large layered graphs further shows that path-task efficiency depends on the particular representation and task \parencite{bae2011quilts}.

However, this research does not directly answer how a game behaviour-tree editor should combine these mechanisms. Behaviour trees differ from ordinary hierarchies such as file directories: sibling order has execution semantics; leaf nodes invoke code; Decorators may prevent a branch from running; and graph state changes on every runtime Tick. A usable solution must preserve execution order, support editing and undo, expose blackboard state, and explain failures without adding a debugging overlay to the game view. Existing research can guide individual design choices, but an integrated, executable and reproducibly evaluated Godot implementation is still required.

\section{Research Aim and Question}

This dissertation aims to design and evaluate composable large-tree display optimisations on a functionally validated Godot 4.6 visual behaviour-tree platform capable of driving real NPCs. The research question is: across different behaviour-tree node scales and physical screen sizes, can the composable display optimisations proposed by this project improve the measurable display and editing-interaction performance of large behaviour trees compared with a traditional baseline display, while preserving behaviour-tree structure and execution semantics?

\section{Research Objectives and Success Criteria}

The project work is divided into platform construction and research evaluation. Platform construction comprises implementing serialisable resources for trees, nodes, Decorators and typed blackboards; implementing the \textsc{success}, \textsc{failure} and \textsc{running} semantics of Root, Sequence, Selector, Reactive Selector, Random Selector, Parallel, Repeat, Wait, Action and Condition; allowing a reusable component to assign a tree and an Actor to an NPC, with leaf nodes calling Actor methods; providing creation, connection, disconnection, dragging, typed editing, validation, saving, loading, undo and redo; and displaying paths, node states, the blackboard and failure reasons only in the Godot editor.

Research evaluation provides independent persistent switches and safe reset paths for large-tree display methods, and uses raw data to compare geometric presentation, target location and context retention across different node scales, physical screen sizes and conditions before and after optimisation. Before conducting the display-optimisation experiments, the fundamental behaviour-tree plugin functions must first be confirmed as correct: each node type must execute as expected; an edited behaviour tree must save and reopen correctly; undo and redo must not damage tree content; and the complex enemy in the playable scene must actually be controlled by the 241-node behaviour tree. Display optimisations must be compared with the original display in a real rendering environment. After an optimisation is disabled, the editor must return to its original display state and must not change the behaviour-tree structure, node positions or execution order. If testing produces an error, crash, memory leak or abnormal exit, the test is judged to have failed even if other checks report success.
